\documentclass[11pt,a4paper]{article}
\usepackage[T1]{fontenc}
\usepackage{lmodern,amsmath,amssymb,booktabs,microtype}
\usepackage[margin=24mm]{geometry}
\usepackage[colorlinks=true,urlcolor=blue,linkcolor=blue]{hyperref}
\newcommand{\EE}{\mathbb E}
\newcommand{\tr}{\operatorname{tr}}
\newcommand{\W}{\mathsf W_2}
\newcommand{\Db}{\overline D}
\title{Central charge windows in the memory spectrum\\A uniform trace bridge and the two force orientations}
\author{Alan Day \and Chengheng\thanks{AI research assistant (OpenAI).}}
\date{Local research note --- 6 November 2026}
\begin{document}\maketitle
\begin{abstract}
Balanced occupation constraints transfer each fixed marked Wick graph
from the full-Fock trace to a compact central charge window with an
$O(p^{-1})$ absolute error in its normalized graph weight. An exact
general-charge Casimir identity then bridges the output deficit to
the two physical force orientations. Existing marked-chord kernels
therefore determine the central oriented spectral matrix and the
complete two-energy measure of the projected second-chaos memory
core. The raw second memory coefficient still requires restoration
of its other descendants at central charge.
\end{abstract}
\section{State and the dynamical result}
Use standard CAR, the full overlapping complex Gaussian $p$-body
Hamiltonian, common couplings in all sectors, and $\hbar=1$.
Let $N$ be even, $\lambda_N=p^2/N\to\lambda>0$,
$m=N/2+xp$, $|x|\le L_0$, and
$C_m=\EE\tr_mH_m^2\to C=J_s^2/\lambda$.
The trace is normalized maximally mixed at charge $m$; all norms
include the Gaussian average. Constants below are uniform for
$x$ in this window and $\lambda_N$ in positive compact sets.

Define $d_i=[H,c_i]$ before restricting neutral operators, and
\begin{equation}
 A_-=\sum_i d_i^\dagger d_i,\quad A_+=\sum_i d_id_i^\dagger,
 \quad O_\pm=\W A_\pm,\quad B=\W h^2,\quad h=H_m.
\end{equation}
Here $\W$ is ordinary coupling-coordinate second chaos.
Put $q=e^{-4\lambda}$, $c=2\lambda(1-q)/(1+q)$ and
$Y_\pm=O_\pm-cB$. The central two-energy limit is
\begin{equation}
 \boxed{\ d\nu_{(Y_-,Y_+)}=
 \begin{pmatrix}w&\eta\\\eta&w\end{pmatrix}d\mu_{q^2},\quad
 \begin{aligned}
 w&=4C^2\lambda+\frac{16C^2\lambda^2q}{1+q},\\
 \eta&=-4C^2\lambda q+\frac{16C^2\lambda^2q}{1+q}.
 \end{aligned}\ }\label{matrix}
\end{equation}
The normalized Mehler measure $\mu_r$ is the existing probe
measure with Hamiltonian variance $C$ and chord parameter $q$.
Its support in each energy is $[-2\sqrt{C/(1-q)},2\sqrt{C/(1-q)}]$.
It is the prior chord kernel, now applied to the present state.
For the output-deficit residual $T$ defined on page 5,
\begin{equation}
 d\nu_T=\frac{C^2v\Delta}{v+\Delta}\,d\mu_{q^2},\quad
 v=2\lambda(1-q)+\frac{16\lambda^2q}{1+q},\quad
 \Delta=2\lambda(1+q).\label{Tmain}
\end{equation}
Both this variance and the eigenvalues of \eqref{matrix} are positive.

\newpage
\section{The occupation ratio on a central window}
For a term fixing $t$ occupied and $L-t$ empty sites, the exact
sector/full-Fock trace ratio is
\begin{equation}
 R_{N,m}(L,t)=2^L\frac{(m)_t^\downarrow(N-m)_{L-t}^\downarrow}
                         {(N)_L^\downarrow}
       =2^L\frac{\binom{N-L}{m-t}}{\binom Nm}.\label{ratio}
\end{equation}
Let $s=L-t$, $d=2t-L$ and $\epsilon=(m-N/2)/N$.
For $L\le kp$ at fixed $k$, all relevant falling factorials are
positive for sufficiently large $p$. An explicit first approximation is
\begin{align}
 A={}&t\log(2m/N)+s\log(2(N-m)/N)\notag\\
 &-\frac{t(t-1)}{2m}-\frac{s(s-1)}{2(N-m)}+\frac{L(L-1)}{2N}.
\end{align}
The elementary log remainder gives
\begin{equation}
 |\log R-A|\le E(m,t)+E(N-m,s)+E(N,L),\quad
 E(a,j)=\frac{j(j-1)(2j-1)}{12a^2[1-(j-1)/a]},\label{error}
\end{equation}
with $E(a,0)=0$. This follows by summing
$|\log(1-y)+y|\le y^2/[2(1-y)]$ for the three factorials.
The denominator alone also gives the global upper bound
\begin{equation}
 \log R\le L\log(1+2|\epsilon|)
                 +\frac{L(L-1)}{2(N-L+1)}\le K_{k,L_0,\lambda}.
 \label{upper}
\end{equation}

Expanding only the inverse denominators in $A$ yields
\begin{align}
 \log R={}&\frac L2\log(1-4\epsilon^2)
       +\frac d2\log\frac{1+2\epsilon}{1-2\epsilon}
       +\frac{L-d^2}{2N}\notag\\
       &+O\left(|\epsilon|L^2/N+L^3/N^2\right).
\end{align}
Since $|\epsilon|\le L_0p/N$, this proves, uniformly in $0\le t\le L$,
\begin{equation}
 |\log R|\le K\left[\frac{1+|d|}{p}+\frac{d^2}{N}\right],\quad
 |R-1|\le e^K|\log R|.\label{balanced}
\end{equation}
The last inequality uses the bounded positive part of $\log R$;
no lower bound on $R$ is needed. A balance estimate on $d$ will
turn this into an integrable error for the actual marked graphs.

\newpage
\section{Transferring fixed marked energy moments}
Fix a Gaussian pairing and a finite number of normal or reverse
commutator marks. Each Gaussian chord carries two $p$-element sets
of opposite orientation. Let $K_{\rm ov}$ be the sum of pairwise
intersection sizes of the distinct sets in this fixed graph.
For a nonzero local occupation term, singleton sites on the two
sets of a chord have opposite occupations and equal original
cardinalities. Removing the multiply used sites changes this
balance by at most a graph-dependent multiple of $K_{\rm ov}$.
Thus
\begin{equation}
 |2t-L|\le c_kK_{\rm ov},\qquad L\le kp.\label{balance}
\end{equation}
For the Wick blocks used here there is no Gaussian chord internal
to a marked block. A force mark, or a fixed number of body-deficit
marks, is supported on intersections and annihilates singleton
sites. Expanding multiply used sites into occupation projectors
has absolute coefficient bound
$(1+K_{\rm ov})^j e^{c_kK_{\rm ov}}$ relative to the common
singleton normalization. Fermionic signs and the orientations are
retained; this is an absolute error estimate on their sum.

The previously used exponential overlap bound can be seen from
the hypergeometric factorial moments of a pair overlap $K_{ij}$:
\begin{equation}
 \EE(K_{ij})_r^\downarrow
 =\frac{(p)_r^{\downarrow,2}}{(N)_r^\downarrow}
 \le\left(\frac{p^2}{N-p}\right)^r.
\end{equation}
Its factorial-moment generating function bounds
$\EE e^{aK_{ij}}$ uniformly for each fixed $a$; H\"older handles
the finite sum of pair overlaps. Combining \eqref{upper},
\eqref{balanced} and \eqref{balance} consequently gives
\begin{equation}
 \left|\mathcal G_{m,N}-\mathcal G_{{\rm Fock},N}\right|
 \le\frac Kp\EE\left[(1+K_{\rm ov})^{j+2}e^{c_kK_{\rm ov}}\right]
 \le\frac{K'}p.\label{graph}
\end{equation}
Here $\mathcal G$ denotes the graph weight with its Gaussian
couplings and singleton powers normalized. Apply this also to
$H^2$ to see that the sector and full-Fock variance normalizations
agree up to $O(p^{-1})$. This does not require summing over all
charge sectors or assuming a tail estimate outside the window.

The oriented-word and Poisson-overlap rules are prior results
\cite{BNR}; their marked versions and endpoint certificates are
already recorded in this project \cite{oriented}. Equation
\eqref{graph} transfers those results to the central trace for
every fixed left/right energy degree. It applies to bounded
coefficient combinations of fixed marked Wick blocks, including
$O_\pm$, $B$ and the previously defined fourth-chaos $Q_4$.
It does not automatically control a raw counterterm multiplied
by a growing factor of $p$, nor a number of energy insertions
growing with $N$.

\newpage
\section{The output deficit and the actual reverse force}
For neutral full-Fock operators define
\begin{equation}
 DX=QX-\sum_i c_i^\dagger Xc_i,\qquad
 \Db X=(N-Q)X-\sum_i c_iXc_i^\dagger.
\end{equation}
At arbitrary charge $m$, CAR gives the exact identities
\begin{equation}
 \mathcal C_xX=(N-m+1)DX+m\Db X-D\Db X,\quad DH=pH,
 \quad (\Db h)_a=\frac{x_a-p(N-m)}{m-p+1}h_a.\label{casimir}
\end{equation}
The compositions are evaluated before restriction. For $u=p-a$ put
\begin{equation}
 D_{\rm in}h=\sum_a uh_a,\quad
 E_{\rm in}=\sum_a\frac{u(u-1+N-2m)}{m-p+1}h_a.
 \qquad (p-\Db)h=2D_{\rm in}h+E_{\rm in}.\label{ein}
\end{equation}
Unlike exact half filling, $E_{\rm in}$ is now $O_{L^r}(p^{-1})$
for every fixed trace/Gaussian moment order needed here. The
factorial input-rank tail proves this from \eqref{ein}; the half-filled
$O(N^{-1})$ bound must not be substituted.

Let $\mathcal S=\sum(a+b-r)\W Z^+_{ab;r}$ and let
$\widetilde{\mathcal S}$ use $(x_{a+b}-x_r)/N$ instead.
Define, before restriction,
\begin{align}
 \overline A&=(2p-\Db)H^2=A_++\{H,(p-\Db)H\},\\
 E_{\rm dbl}&=(2p-D)\overline A,\qquad
 U_{{\rm in},2}=\{h,D_{\rm in}^2h\}+2(D_{\rm in}h)^2.
\end{align}
Substituting into \eqref{casimir} gives the general-charge bridge
\begin{align}
 N\left(\widetilde{\mathcal S}-\frac{O_-+O_+}{2}\right)
 ={}&(N/2-m+1-2p)O_-+(m-N/2-2p)O_+\notag\\
 &+(2m-N-1)\W\{h,D_{\rm in}h\}-\W U_{{\rm in},2}
   +\W E_{\rm dbl}\notag\\
 &+(m-2p)\W\{h,E_{\rm in}\}.\label{bridge}
\end{align}
The double mark is supported on overlap sites: each deficit
annihilates a singleton and a preceding deficit does not create
a new singleton contribution. Its coefficient bound is
$K(1+K_{\rm ov})^2$, so \eqref{graph} bounds its fixed moments.
All other terms are controlled by the established input tails
and Gaussian product bounds. Divide \eqref{bridge} by $N$ and
use $m-N/2=O(p)$. The positive output moments also give
$\|\mathcal S-\widetilde{\mathcal S}\|_2=O(p^{-1})$. Hence
\begin{equation}
 \boxed{\quad\mathcal S=(O_-+O_+)/2+O_{L^2}(p^{-1})\quad}
 \label{strong}
\end{equation}
uniformly on the central window, with bounded physical $C_m$.

\newpage
\section{The spectral measure of the projected quadratic core}
The transferred force graph calculation gives \eqref{matrix} for
all fixed mixed energy moments. The existing $Q_4$ endpoint
certificates also give $\nu_{Q_4,Y_-}=\nu_{Q_4,Y_+}=0$ centrally.
These are entire cross measures: finite-system matrix measures
\begin{equation}
 \nu_{ij,N}(dE,dE')=\EE\tr_m[X_i\Pi_h(dE)X_j\Pi_h(dE')]
\end{equation}
are positive. Uniform diagonal moments give tightness and uniform
integrability, and the candidate compactly supported measures
are moment determinate. Spectral Cauchy--Schwarz controls the
off-diagonal measures. A subsequence argument gives convergence
uniformly over the stated central parameters.

Let $\Pi_N$ be the actual finite-$N$ orthogonal projector onto
$\operatorname{span}(B,L_{\rm in},O_-)$ and
$T_N=(1-\Pi_N)\mathcal S$. The previous complete central Gram
has a nondegenerate limiting inverse. With $\rho=v/(v+\Delta)$,
\begin{equation}
 T_N=\left(\frac12-\rho\right)Y_-+\frac12Y_++o_{L^2}(1).
\end{equation}
The $L_{\rm in}$ projection coefficient tends to zero. Multiplying
the $2\times2$ matrix in \eqref{matrix} proves \eqref{Tmain} and
\begin{equation}
 \nu_{T,Y_-}=0,\qquad \nu_{T,Q_4}=0.
\end{equation}
For an $L^2$ replacement the difference of self measures is bounded
in total variation by
$\|X-X'\|_2(\|X\|_2+\|X'\|_2)$. Thus \eqref{strong} and the
projection replacement transfer the weak spectral limit and all
bounded spectral tests, without expanding a time exponential.
Unbounded polynomial tests on a replaced field require their own
uniform integrability; the marked fields have it from \eqref{graph}.

Recall the actual raw core
\begin{equation}
 \mathcal E=\sum_i[H,[H,c_i]]^\dagger[H,[H,c_i]]+[h,[h,A_-]]/4.
\end{equation}
The established central strong expansion implies
$\mathcal R_{2,N}:=(1-\Pi_N)\W\mathcal E
=8\lambda_NC_mq_N^2T_N+o_{L^2}(1)$. Therefore
\begin{equation}
 \boxed{\ d\nu_{\mathcal R_2}=V_2\,d\mu_{q^2},\quad
 V_2=64\lambda^2C^4q^4\frac{v\Delta}{v+\Delta},\quad
 C_{\mathcal R_2}(t)=V_2G_{q^2}(t).\ }\label{core}
\end{equation}
The limiting measure has no static atom at fixed $0<q<1$.
Its normalized second frequency moment is $2C(1-q^2)$.
With $C=J_s^2/\lambda$, after DS and then $\lambda\downarrow0$,
$V_2\to1024J_s^8$ and $G_{q^2}(t)\to\operatorname{sech}^4(2J_st)$.
Equation \eqref{core} identifies the projected second-chaos core;
it is not yet a restoration of the complete central raw $U_2$.

\newpage
\section{Executed checks and continuation}
Three new scientific records are stored in
\texttt{results/susskind-memory-central-spectrum}:
\begin{itemize}
 \item \texttt{trace-ratio.json}: 24 exact finite binomial ratios
 and 324 central cases, with $\lambda=0.1,0.25,0.5$,
 $p=16,32,64,128$, $x=-2,0,2$, $L/p=2,4,8$, and imbalance
 $2t-L=0,2,6$. Every case satisfies the explicit log remainder
 \eqref{error} and upper bound \eqref{upper}. The 90/130-digit
 refinement differs by at most $2.23\,10^{-85}$.
 \item \texttt{force-bridge.json}: 108 exact cleared-rational
 identities on six non-half-filled $(N,p,m)$ systems and three
 integer Hermitian inputs each. The checks use full CAR and
 complete Casimir projectors, including the general $E_{\rm in}$
 correction and the pre-restriction compositions in \eqref{casimir}.
 \item \texttt{spectrum-assembly.json}: symbolic oriented matrix,
 residual projection, zero cross with $Y_-$, and core normalization.
 It verifies the hashes of the prior local-mark and endpoint
 certificates rather than rerunning them. The reverse certificate
 retains 128 templates and 6,072 words; its coefficients vanish
 for arbitrary fixed energy degrees by the original endpoint proof.
\end{itemize}
For $\lambda=0.5$, $x=2$, $L=4p$ and $2t-L=2$, the actual trace
ratios at $p=16,32,64,128$ are approximately
$0.7985081$, $0.9027791$, $0.9522632$, $0.9763479$.
The corresponding $p(R-1)$ values approach a finite number near
$-3$ in this example. Finite-size differences are retained.
The analytic remainder bounds are not interval-rounding certificates;
the numerical evaluations use precision refinement.

The next missing pieces for the raw central $U_2$ are its lower-chaos
descendant table and exact conserved subtraction, and the filtered
linear-field interfaces needed in its mixed spectral matrix.
The new trace theorem handles fixed marked blocks but must not
silently be multiplied by a divergent counterterm. These tasks
precede the charge tails, finite-memory two-insertion remainder,
and Susskind's bulk meson space/mass reconstruction.
No bound uniform in growing word length, growing $|x|$, or joint
small $\lambda$ and finite $N$ is asserted.

\begin{thebibliography}{9}\small
\bibitem{BNR} M. Berkooz, V. Narovlansky and H. Raj,
\emph{Complex Sachdev--Ye--Kitaev model in the double scaling limit},
\href{https://arxiv.org/html/2006.13983v2}{arXiv:2006.13983v2},
Sec. 3.3, Eqs. (3.11)--(3.15). Existing oriented trace and overlap rules.
\bibitem{oriented} A. Day and Chengheng, local project notes
\emph{Two orientations of the fermionic force},
\emph{Central charge windows}, and
\emph{Closing the central-window memory Gram} (6 November 2026).
Their formulas and executed certificates are inputs to this extension.
\end{thebibliography}
\end{document}
